\subsection{\texorpdfstring{SERIES IN \(\mathbf{R}^n\)}{SERIES IN R TO THE N}}

A series whose general term is \(\boldsymbol{x}_{m} = (x_{mi})_{1 \leqslant i \leqslant n}\) converges in \(R^{n}\) if and only if each of the n series \((x_{mi})_{m \in \mathbb{N}}\) converges in R.

DEFINITION 1. A series of points of \(\mathbf{R}^n\) is said to be absolutely convergent if the series of Euclidean norms of its terms is convergent.

PROPOSITION 4. A series of points of \(\mathbf{R}^n\) is commutatively convergent if and only if it is absolutely convergent.

This is a consequence of Proposition 9 of Chapter III, § 5, no. 7 and Theorem 1 above.

The examples given in Chapter IV, § 7 show that a series in \(\mathbf{R}^n\) can be convergent without being absolutely convergent.
% semantic-fragments: legacy-input-seam
\relax{}
\ExerciseGroupsStart
